\documentclass[11pt,letterpaper]{article}
\usepackage[T1]{fontenc}
\usepackage{lmodern,amsmath,amssymb,amsthm,microtype}
\usepackage[margin=1in]{geometry}
\usepackage[hidelinks]{hyperref}
\providecommand{\tightlist}{\setlength{\itemsep}{0pt}\setlength{\parskip}{0pt}}
\DeclareUnicodeCharacter{2264}{\ensuremath{\le}}
\DeclareUnicodeCharacter{2265}{\ensuremath{\ge}}
\DeclareUnicodeCharacter{2260}{\ensuremath{\ne}}
\DeclareUnicodeCharacter{2212}{\ensuremath{-}}

\DeclareUnicodeCharacter{2211}{\ensuremath{\sum}}

\usepackage{longtable,booktabs,array}
\title{Orbit packing rules out eighth weights below 132 in the nine-weight search}
\author{Aafi Khan Malek\\ \small Independent researcher, Ahmedabad, India\\
\small\href{mailto:aafimalek2032@gmail.com}{aafimalek2032@gmail.com}}
\date{28 November 2026}
\begin{document}
\maketitle
\textbf{Theorem.} Let nine sorted positive real weights have pairwise
unit-separated subset sums. If \(149\le x_9<150\), then \(x_8\ge132\).
The proof extends the
\href{https://github.com/Aafimalek/novel-discoveries/tree/484f4f670bba27fa3c0eec23ceae0db0a923940a/nine-separated-subset-sums/verification/rank-cuts-below-150/penultimate-cap-130}{previous rank-gap argument}
by assigning the sign orbits to three regions inside every residual
deficit cell. It does not determine the optimum of the nine-weight
problem and has not been checked in Lean.

The forced rank starts are 142 through 172, with slots
\(q=34,\dots,64\). The common counts are \(64\le a\le74\) and
\(b=148-a\). Each \(\alpha_q=u_{a-q+1}\) lies in \([149-2q,151-2q)\).
The resulting bins for the 128 complementary signs are:

\begin{longtable}[]{@{}lr@{}}
\toprule\noalign{}
Half-open range & Capacity \\
\midrule\noalign{}
\endhead
\bottomrule\noalign{}
\endlastfoot
{[}1,21) & \(a-64\) \\
each {[}149−2q,151−2q), q from 64 down to 34 & 1 \\
{[}83,217) & 33 \\
{[}278,321) & \(148-2a\) \\
{[}381,10000) & \(a-20\) \\
\end{longtable}

The two wide gaps have width at least
\(\alpha_{34}-\alpha_{64}+4\ge64\). The total deficit relative to
\(y=(125,\ldots,132)\) is \(R=\sum_i(y_i-x_i)\). If a prescribed form is
unbalanced, its cap magnitude is at least 242 rather than below 83, so
\(R>159\). Otherwise \(\alpha_{34}\ge81\) is balanced, and its cap
magnitude is at most 16, so \(R\ge65\). Thus \(R\ge65\) in either case.

\section{Sign orbit packing in deficit
cells}\label{sign-orbit-packing-in-deficit-cells}

Write \(W=32\), half the certified wide-gap width. For a base coordinate
strictly smaller than \(W\), flipping its sign changes absolute
magnitude by less than \(2W\); the same holds for swapping signs at
adjacent weights whose difference is strictly smaller than \(W\). Either
operation therefore preserves each set of signs below either wide gap.
There is a nonsmall coordinate since otherwise all eight-coordinate
magnitudes would be below \(8W\), contrary to the last bin's lower bound
and positive capacity. Representing complementary pairs by fixing that
coordinate's sign shows that if \(s\) coordinates are smaller than
\(W\), then \(2^s\) divides \(a,b\) and hence 148. Therefore \(s\le2\),
\(x_3\ge32\), and the ordered deficits cap the third deficit at
\(127-32=95\). On cells which certify one or two small coordinates, use
the implied divisibility too.

The exact checker bounds each fixed deficit by its interval endpoints
and bounds each remaining ordered deficit separately. If the remaining
tail has \(N\) entries capped at \(c\) and total in \([L,U]\), its
\(j\)-th entry is at least \((L-(j-1)c)/(N-j+1)\) and at most \(U/j\).
Safe outward rounding gives estimates sufficient to certify some flips
and adjacent swaps throughout the whole cell. An adjacent gap is
\(1+r_i-r_{i+1}\), so its upper endpoint must be \textbf{strictly} below
\(W\) before adding the generator. The generator action on complementary
sign pairs is computed by restoring the first sign to plus when
necessary.

The orbits of the certified generators must each belong entirely to one
of three broad regions: ranks at most \(a\), strictly between counts
\(a\) and \(b\), or above count \(b\). Their exact capacities are
\(a,b-a,128-b\). A two-dimensional bitset dynamic program asks whether
unrestricted orbit sizes can fill these capacities. If so, construct the
usual outward absolute magnitude interval for each sign (fixed deficit
boxes and prefix-sign averages on the ordered tail). Each orbit may take
only broad region labels that are possible for \textbf{every} sign in
that orbit according to these intervals. The same exact bitset
recurrence with the allowed labels checks this necessary condition.
Ordinary full-bin interval matching is checked independently. This
deliberately uses relaxations; rejecting any one is sufficient, and
passing one is not deemed feasible.


\section{Exact replay}
The two-deficit matching steps leave 7,328 profiles. Sign-flip divisibility
rejects 299 before further splitting. The remaining profiles give 279,322
three-deficit cells. Tail and moment constraints reject 624. Orbit sizes
alone reject 139,190 more, ordinary interval conditions 40,707, and the
joint orbit-membership conditions 98,801. None survives.
Use \texttt{python3 caps\_orbit.py 131} for the standard-Python replay.
\end{document}
